\section{K-034: rank-two odd inertia is confined to the conductor downstairs}

The exact normalization series of K-033 is insufficient by itself to
exclude higher odd grading stabilizers.  The following support argument
does exclude them everywhere on that normalization.  It applies to an
entire class of rank-two block models and gives a coherent-class
reduction to their conductors, without assuming that the component is
normal.

\subsection*{A small-denominator lemma for the block model}

Let $\Sigma\subseteq\{1,\ldots,m\}^3$ be free, with linearly
independent incidence columns and a positive margin-one mass vector
$(p_e)_{e\in\Sigma}$.  Fix an edge $e_0$.  Assume that every $p_e$
has denominator dividing a power of two $M$.

Replace its selected coordinate in each factor by a two-dimensional
block.  Let $U\subseteq(\mathbb C^{m+1})^{\otimes3}$ be the complete
expanded support space: $e_0$ expands to a $2\times2\times2$ block;
an edge sharing one selected coordinate expands to a two-vector;
and the remaining edges are scalar coordinates.  Freeness excludes
an edge sharing exactly two selected coordinates.
Write $L$ for the block-diagonal group
\[
                  L=(GL_2\times(\mathbb C^*)^{m-1})^3,
\]
and $C=L\cap SL_{m+1}^3$.

\begin{lemma}[K-034a]
If $v\in U$ is special-linear semistable, then every
$l\in\operatorname{Stab}_L(v)$ satisfies
\[
                             \chi(l)^{2M}=1.           \tag{K-034a}
\]
Consequently, whenever $U$ is a complete fixed space as in
K-032 or K-033, every nonzero point of $U//C$ has original
grading stabilizer contained in $\mu_{2M}$.
\end{lemma}

\begin{proof}
We first prove a statement about any support in $U$ admitting
nonnegative margin-one masses.  Collapse the two-dimensional blocks
to the selected old coordinates.  Aggregate the masses in each old
edge group.  Their old incidence margins are one outside the three
selected coordinates and two at those selected coordinates.  Since
the old incidence columns are independent, the aggregate masses
are forced to be
\[
             p_e\quad(e\ne e_0),\qquad p_{e_0}+1\quad(e=e_0).
                                                        \tag{K-034b}
\]
All have denominator dividing $M$.

Consider the bounded nonempty polytope of nonnegative masses with
these group totals and with first-coordinate margin one in each
of the three two-dimensional blocks.  These conditions imply all
the original margins: the second block margin is two minus the
first, and the outside margins are fixed by (K-034b).  Choose a
vertex.  In each group choose one positive variable as a base
variable, and eliminate it using the group total.  The remaining
positive variables have linearly independent columns in the three
remaining equations.  Otherwise a small perturbation of both signs
would preserve all constraints and positivity, contradicting the
vertex condition.  In particular there are at most three of them.

For the tensor group, the columns after this elimination are
differences $a-a_0$ of vertices of the unit cube $\{0,1\}^3$, with
a common base vertex $a_0$.  Change the sign of row $i$ when the
$i$th coordinate of $a_0$ is one.  All these columns then become
binary columns.  A vector group contributes only a signed unit
column, which becomes a binary unit column after a column sign
change.  Thus the absolute values of all relevant minors are
those of binary matrices with at most three rows and columns.

A nonsingular such minor has absolute value one or two.  For the
only nontrivial case, a binary $3\times3$ matrix $D$, form the
$4\times4$ sign matrix
\[
                 \begin{pmatrix}1&\mathbf1^t\\
                       \mathbf1&\mathbf1\mathbf1^t-2D
                       \end{pmatrix}.
\]
Its determinant is $(-2)^3\det D$, while Hadamard's inequality
bounds its absolute value by $16$.  Hence $|\det D|\le2$.
The smaller minors are immediate.  Cramer's rule in an independent
square row minor now shows that all the remaining variables,
and therefore all the eliminated base variables, have denominators
dividing $2M$.  We have proved that every feasible support in $U$
admits a margin-one mass vector with this common denominator.

Return to the stabilizer element $l$.  Its semisimple and unipotent
parts lie in the algebraic stabilizer; the unipotent part has
determinant character one.  Conjugate its semisimple part within
$L$ to a diagonal element.  This conjugation preserves $U$ and
special-linear semistability.  Torus semistability supplies
nonnegative margin-one masses on the resulting tensor support.
Use the mass vector $q_e$ of denominator dividing $2M$ just proved.
If $\lambda_{a,i}$ are the diagonal eigenvalues, then
\[
 \chi(l)^{2M}
  =\prod_{a,i}\lambda_{a,i}^{2M}
  =\prod_{e=(i,j,k)}
       (\lambda_{1,i}\lambda_{2,j}\lambda_{3,k})^{2M q_e}=1,
\]
because every supported tensor coordinate is fixed.  This proves
(K-034a), including for a semisimple element of arbitrary order.

Finally let $y\ne0$ be a point of $U//C$.  Choose its closed
$C$-orbit representative.  Equality of nullcones, proved in K-032
by the unitary moment-map argument, makes it special-linear
semistable.  The action of $L$ on this quotient factors through
the determinant-product character, with the original tensor
grading.  If a grading element $\eta$ fixes $y$, choose $l_0\in L$
with $\chi(l_0)=\eta$.  The vectors $l_0v$ and $v$ have the same
quotient and closed $C$-orbits, so $l_0v=cv$ for some $c\in C$.
Then $c^{-1}l_0$ fixes $v$ and has character $\eta$.  Applying
(K-034a) gives $\eta^{2M}=1$, as required.
\end{proof}

\subsection*{All coherent classes upstairs are three-integral}

Apply the lemma to the actual normalization $\nu:\widetilde Y\to Y$
in K-033, with $M=512$.  The kernel $\mu_M$ acts trivially on this
normalization, and every nonzero grading stabilizer is either
$\mu_M$ or $\mu_{2M}$.  In the regrading $z=t^M$, its stabilizers
therefore have order one or two.  In particular
\[
              (\widetilde Y^{\mu_{3^a}})_{\rm red}=\{0\}
                       \qquad(a\ge1)                  \tag{K-034c}
\]
in the regraded convention.

It follows that every finitely generated graded module over
$B=\mathbb C[\widetilde Y]$ has a three-integral Laurent expansion
at a primitive $M$th root in the original grading.  Here the
assertion concerns all modules, not only $B$ itself.  Indeed,
the finite-length Koszul argument of K-027 excludes poles at
every order $M3^a$, $a\ge1$, for such a module.  The same
cyclotomic no-collision lemma then gives integrality of every
Laurent coefficient at $\xi$, including the negative coefficients.
Arbitrary finite graded shifts are allowed, exactly as in K-027.

More explicitly, an arbitrary graded module decomposes by degrees
modulo $M$, so its Hilbert series is
$\sum_{a=0}^{M-1}t^aF_a(t^M)$.  Each $F_a(z)$ has reduced
denominator consisting only of powers of $1-z$ and $1+z$:
the same Koszul argument excludes every other root of unity,
since no nonzero point has any other finite grading stabilizer.
This universal assertion does not follow from the particular
Hilbert numerator calculated in K-033; it follows from K-034a.

\subsection*{A conductor-supported refinement of the actual exceptional class}

Let $Y_0,Y_1,Y_2$ be the cyclic orbit of the actual component $Y$,
and let $\nu_i:\widetilde Y_i\to Y_i$ be their normalizations.
Let $C_i\subseteq Y_i$ be the conductor locus, the support of
$\nu_{i*}\mathcal O_{\widetilde Y_i}/\mathcal O_{Y_i}$.
Let $D_i$ be the union of $C_i$ with the intersections of $Y_i$
with the other reduced irreducible components of $X_{10}^{\mu_M}$.
These are closed subsets of dimension at most three.

\begin{proposition}[K-034b]
In the integral coherent-class reduction K-031, this cyclic orbit
of four-dimensional components can be replaced in the exceptional
support by $D_0\cup D_1\cup D_2$.  This replacement is valid for
the entire Laurent principal part.  It leaves every other
exceptional component unchanged and makes no assertion that the
remaining conductor-supported class is zero or three-integral.
\end{proposition}

\begin{proof}
Work first with a graded coherent module $F$ on one $Y_i$.
The natural homomorphism
\[
 F\longrightarrow \nu_{i*}
             (\mathcal O_{\widetilde Y_i}\otimes_{\mathcal O_{Y_i}}F)
\]
is an isomorphism away from $C_i$.  Its kernel and cokernel are
coherent and supported on $C_i$.  Thus in integral graded $G_0$
\[
                [F]=\nu_{i*}[\widetilde F]+\delta_i,
             \qquad \operatorname{supp}\delta_i\subseteq C_i,
                                                        \tag{K-034d}
\]
where $\widetilde F$ is the ordinary pullback module.  No flatness,
derived pullback, or division by a finite degree is needed.

Use coherent localization away from the component intersections
to separate the original class along the three open component
pieces.  Choose an extension on $Y_0$ and transport it to $Y_1$
and $Y_2$; the remainder is supported on the intersections.
Nilpotent structures are handled by the same coherent
d\'evissage used in K-027--K-031.  Apply (K-034d) to each transported
class.  The three upstairs classes have identical Hilbert series.
Each such series is three-integral by (K-034c) and the preceding
Koszul argument.  Their sum is therefore three times a
three-integral Laurent series.  After division by the moving
generator factor $D_M(t)$ of K-027, this remains true, because
$D_M(\xi)$ is a three-adic unit.

Consequently the entire contribution of these upstairs classes
vanishes modulo three in every Laurent coefficient.  All
remaining classes associated with this orbit are supported on
the conductors or component intersections, precisely the sets
$D_i$.  Their dimension is at most three: normalization is
birational and an isomorphism on a dense open set, and distinct
irreducible components meet in proper closed subsets.  This
proves the claimed refinement of the support in K-031.
\end{proof}

This result applies whether or not the downstairs component
actually contains a higher odd-inertia point; that incidence has
not been classified here.  If such a point exists on this
component, K-034c forces it into the conductor, with its
normalization branches permuted.  The proposition is therefore a
reduction of the genuine coherent-class obstruction, rather than
an inference of component normality from its Hilbert series.
It does not compute the conductor class and does not settle the
global exact-$M$ pole.

\paragraph{Method and attribution.}
The binary-minor estimate is an elementary application of the
classical Hadamard determinant inequality, not a new general
determinant bound.  The finite-normalization conductor sequence
and integral coherent localization are standard methods; the
exact kernel--cokernel construction used here is displayed in
(K-034d), and the required finite-length Koszul and cyclotomic
arguments are the already proved lemmas of K-027.  The contribution
of this result is their application to the specified tensor block
models and the resulting support reduction.  No historical
priority claim is made for the underlying methods.
